\chapter{Power spectral density and Welch's method}
\label{app:psd}

The spectra in Chapter~\ref{ch:dataset} compare fluctuations in recorded altitude.
They contain both flight dynamics and measurement error. A spectral estimate does
not separate those contributions without further assumptions.

\section{From covariance to a spectrum}

Let $x_n=z_n-\mathbb E[z_n]$ be a real, uniformly sampled, second-order stationary
process, with sampling interval $\Delta t=1/f_s$ and covariance
$C_k=\mathbb E[x_nx_{n+k}]$. When the covariance is summable, its two-sided spectral
density is
\begin{equation}
 S_2(f)=\Delta t\sum_{k=-\infty}^{\infty}C_k
              e^{-2\pi i f k\Delta t},\qquad |f|\leq f_s/2.
\end{equation}
For real signals $S_2$ is even. The one-sided density doubles the positive-frequency
contribution, giving
\begin{equation}
 \int_0^{f_s/2} S_1(f)\,df=C_0=\operatorname{Var}(x_n).
 \label{eq:psd-variance}
\end{equation}
Its units are $\mathrm{m^2\,Hz^{-1}}$ for an altitude signal. The Nyquist frequency
$f_s/2$ limits distinguishable sampled frequencies: physical motion above that
frequency may be aliased into the observed band.

A full flight altitude series need not be stationary. Equation~\eqref{eq:psd-variance}
therefore gives the stationary reference interpretation; it is not a claim that
the variance of an entire climbing and descending flight equals the integral of
its detrended Welch estimate. Here the spectrum is a finite-record diagnostic,
computed after detrending shorter windows.

\section{The estimator used here}

For window $j$ of length $L$, subtract its least-squares linear trend and denote the
residuals by $y_{j,n}$. With a periodic Hann taper $w_n$, define
\begin{equation}
 I_j(f_k)=\frac{\left|\sum_{n=0}^{L-1}w_n y_{j,n}
                    e^{-2\pi i kn/L}\right|^2}
                   {f_s\sum_{n=0}^{L-1}w_n^2},\qquad f_k=kf_s/L.
\end{equation}
This is the two-sided normalisation. For a one-sided estimate, positive bins are
doubled except for the Nyquist bin when $L$ is even; the zero-frequency bin is
also left unchanged. Welch's estimate averages the resulting periodograms over
$K$ overlapping windows~\cite{welch1967}:
\begin{equation}
 \widehat S_1(f_k)=\frac1K\sum_{j=1}^{K}I_{j,1}(f_k).
 \label{eq:welch-estimator}
\end{equation}
The implementation uses $L=256$, half-window overlap, linear detrending and the
arithmetic mean within each flight. The frequency spacing is $f_s/L$; tapering
broadens the effective resolution. Details of the eligible intervals and the
separate aggregation across flights are given in \nameref{sec:altchannel}.

An individual periodogram has substantial sampling variability. Under suitable
stationarity and weak-dependence conditions, averaging independent windows reduces
its variance approximately by $K$. Overlapping windows are correlated, so their
effective number is smaller. Increasing $L$ improves frequency resolution but
leaves fewer windows to average; keeping $L$ fixed also retains smoothing bias
as the record grows. The mean Welch estimator is not, by itself, resistant to
arbitrarily large outliers.

Linear detrending suppresses slow variation within each window; it does not remove
all physical motion or isolate sensor noise. Likewise, a high-frequency plateau
can be compatible with measurement noise without uniquely identifying its cause.
Comparisons between GNSS and barometric spectra must account for missing-channel
selection, sampling, interpolation and the response of any subsequent filter.
